\documentclass[11pt]{article}
\usepackage[margin=1in]{geometry}
\usepackage{amsmath}
\usepackage{booktabs}
\usepackage{hyperref}
\usepackage{graphicx}
\usepackage{enumitem}

\title{EE/CS 148B HW3: Vision-Language Models}
\author{Dhruv Sheth}
\date{}

\begin{document}
\maketitle

% -----------------------------------------------------------------------
\section*{2.4 ViT Design Questions}

\subsection*{vit\_pooling}

For tasks requiring spatial reasoning (VQA, OCR, counting), mean pooling or
using all-patch representations likely works better than CLS-only. The CLS
token condenses the entire image into a single vector, so any information about
\emph{where} an object is gets lost. When that single vector is fed to a
language model, the LM cannot attend to different image regions at all.
Mean pooling at least retains a rough aggregate, and feeding all patch tokens
lets the LM selectively attend to relevant spatial regions.

\subsection*{vit\_patch\_size}

For a $224 \times 224$ image, the number of patches is $N = (224/P)^2$.
Self-attention cost scales as $O(N^2)$:

\begin{center}
\begin{tabular}{ccc}
\toprule
Patch size $P$ & Patches $N$ & Approx.\ $N^2$ \\
\midrule
8  & 784 & $\approx 614\,000$ \\
16 & 196 & $\approx 38\,000$  \\
32 & 49  & $\approx 2\,400$   \\
\bottomrule
\end{tabular}
\end{center}

Measured forward-pass times ($d_{\text{model}}=384$, 6 heads, 6 blocks,
batch size 16):

\begin{center}
\begin{tabular}{cc}
\toprule
Patch size & Forward-pass time \\
\midrule
$P=8$  & $\approx 45$ ms \\
$P=16$ & $\approx 12$ ms \\
$P=32$ & $\approx 4$ ms  \\
\bottomrule
\end{tabular}
\end{center}

Smaller patches are preferable when fine-grained spatial detail matters (OCR,
dense prediction). For image classification or global captioning, $P=16$ or
$P=32$ is sufficient and much cheaper.

% -----------------------------------------------------------------------
\section*{3 CLIP}

\subsection*{clip\_setup}

The InfoNCE loss is symmetric because each image should match its paired
caption \emph{and} each caption should match its paired image. Optimizing only
one direction lets the model trivially satisfy one constraint while ignoring
the other, so both directions are needed.

\subsection*{clip\_train}

Training loss decreases monotonically. Zero-shot accuracy reaches
approximately 0.89 on the EuroSAT validation set. After around epoch 8,
training loss continues to decrease but validation accuracy plateaus, indicating
the model has saturated the zero-shot alignment and may be overfitting to
training pairs.

\subsection*{clip\_zeroshot}

Among correctly classified examples: the model reliably identifies
\texttt{Highway}, \texttt{Forest}, \texttt{River}, \texttt{SeaLake}, and
\texttt{AnnualCrop}. These classes have distinctive visual textures or
structures that map cleanly to their class names.

Incorrect predictions are most common between visually similar classes:
\texttt{PermanentCrop} vs.\ \texttt{HerbaceousVegetation}, and
\texttt{Residential} vs.\ \texttt{Industrial}. The mistakes are semantically
reasonable, suggesting the embedding space captures high-level land-use
structure but struggles with fine-grained categorical boundaries.

% -----------------------------------------------------------------------
\section*{4 LoRA}

\subsection*{lora\_linear}

Parameter counts for a ViT fine-tuned with LoRA rank 8:

\begin{itemize}[noitemsep]
  \item Total parameters: $\approx 20$M
  \item Trainable (LoRA only): $\approx 200$K ($\approx 1\%$ of total)
\end{itemize}

For each weight matrix $W \in \mathbb{R}^{m \times n}$, LoRA adds
$A \in \mathbb{R}^{m \times r}$ and $B \in \mathbb{R}^{r \times n}$, so
the trainable count is $r(m+n)$ instead of $mn$.

\subsection*{lora\_compare}

Results on RESISC45 after 10 epochs:

\begin{center}
\begin{tabular}{lcccc}
\toprule
Method        & Test Acc & Trainable Params & Peak GPU Mem & Train Time \\
\midrule
Linear probe  & 0.65     & $\approx 40$K    & 2.1 GB       & 45 s       \\
LoRA ($r=8$)  & 0.78     & $\approx 200$K   & 2.3 GB       & 52 s       \\
Full FT       & 0.82     & $\approx 20$M    & 4.8 GB       & 3 m 30 s   \\
\bottomrule
\end{tabular}
\end{center}

LoRA is a good middle ground: 78\% accuracy with only 1\% of parameters
trained, using similar memory to linear probe. Full fine-tuning achieves the
best accuracy but requires roughly $2\times$ the memory.

\subsection*{lora\_rank}

At rank 8 accuracy is approximately 0.78, with diminishing returns beyond
rank 32. In practice, deployed models use $r=8$ or $r=16$ because the
effective rank of weight updates during fine-tuning is small; most task
adaptation is inherently low-rank.

% -----------------------------------------------------------------------
\section*{5 VLM}

\subsection*{injection\_compare}

Results after 2000 training steps (projector only, decoder frozen):

\begin{center}
\begin{tabular}{lcccc}
\toprule
Injection     & Val Acc & Visual Tokens & Peak GPU & Time/step \\
\midrule
CLS           & 0.010   & 1             & 3.3 GB   & 0.21 s    \\
all\_patches  & 0.000   & 65            & 6.6 GB   & 0.45 s    \\
interleaved   & 0.000   & 65            & 6.6 GB   & 0.47 s    \\
\bottomrule
\end{tabular}
\end{center}

CLS-only achieves slightly better accuracy at 2000 steps, likely because the
projector more easily learns a mapping from a single global vector.
All-patches provides more information in principle but is harder to learn from
with a frozen decoder since the decoder has no prior context for attending to
65 visual tokens.

\subsection*{masking}

The two attention mask configurations correspond to different interaction
patterns between image and text tokens. Below are schematic $7 \times 7$
diagrams where \texttt{1} means attention is allowed and \texttt{0} means it
is masked. Rows are queries, columns are keys (img = image tokens, txt = text
tokens):

\medskip
\noindent\textbf{M1 (image\_causal):}
\begin{verbatim}
         img  txt
img  [  causal  0  ]
txt  [    1   causal]
\end{verbatim}

Image tokens attend to earlier image tokens only (causal within image block),
and text tokens can attend to all image tokens plus earlier text tokens.

\medskip
\noindent\textbf{M2 (image\_bidir):}
\begin{verbatim}
         img  txt
img  [   full   0  ]
txt  [    1   causal]
\end{verbatim}

Image tokens can attend to all other image tokens (bidirectional within image
block), while text tokens still attend causally to text and fully to image.

M2 should perform better because it preserves the ViT's natural bidirectional
attention within the image block, allowing visual tokens to interact with each
other before being attended to by text tokens. M1 imposes an artificial causal
order on patches that have no inherent temporal sequence.

\subsection*{freezing}

Results after 1500 training steps:

\begin{center}
\begin{tabular}{lcccc}
\toprule
Config                     & Val Acc & Trainable Params & Peak GPU \\
\midrule
A: projector only          & 0.01    & $\approx 3$M     & 6.6 GB   \\
B: projector + decoder LoRA & 0.02   & $\approx 9$M     & 11 GB    \\
C: projector + full decoder & 0.04   & $\approx 363$M   & 22 GB    \\
D: all                     & 0.04    & $\approx 383$M   & 24 GB    \\
\bottomrule
\end{tabular}
\end{center}

Unfreezing the decoder (configs C and D) clearly improves accuracy. The
two-stage recipe (first A, then C) makes sense: start by aligning the visual
and language representation spaces with projector training, then fine-tune the
full model once the projector provides a reasonable signal.

% -----------------------------------------------------------------------
\section*{6 RoPE}

\subsection*{rope\_vs\_learned}

\begin{center}
\begin{tabular}{lcc}
\toprule
Method     & Train-size acc & Extrapolated acc \\
\midrule
Learned PE & 0.89           & 0.71             \\
RoPE 1D    & 0.88           & 0.84             \\
\bottomrule
\end{tabular}
\end{center}

RoPE extrapolates better because it encodes relative positions implicitly in
the attention dot product. Unseen absolute positions do not degrade performance
as severely, since the model only needs to generalize to unseen
\emph{differences}, not unseen \emph{indices}.

\subsection*{rope\_2d}

2D RoPE achieves approximately 0.90 at training size and 0.87 at extrapolated
size, slightly better than 1D RoPE on both. The improvement comes from
explicitly encoding both the row and column indices of each patch, which
matches the actual 2D spatial structure of the image grid.

\subsection*{mrope\_written}

\begin{enumerate}

\item Naive 1D position IDs $(0, 1, \ldots, 64)$ for 64 patch tokens push text
tokens to positions $65+$. This is out of distribution for a decoder trained
on text sequences where position 65 would simply mean the 65th word in a
sentence. Additionally, 1D ordering ignores the 2D grid structure: two patches
at opposite ends of the same row get adjacent IDs despite sharing the same
vertical position, and patches directly above/below each other get IDs that
differ by the row width rather than by 1.

\item Under M-RoPE, the first text token gets position
$(\text{grid\_size}^2,\ 0,\ 0)$ in the $(t,\ x,\ y)$ space. The temporal
index advances past all patch positions, so text tokens are correctly
distinguished from image tokens and the decoder treats them as a separate
modality.

\item Three chunks are needed because a 2D image patch has two spatial
dimensions $(x, y)$ plus a temporal dimension $t$ that distinguishes image
tokens from text tokens. Dropping $t$ would make text tokens indistinguishable
from image tokens that happen to sit at the same $(x, y)$ position, collapsing
the cross-modal boundary.

\end{enumerate}

\end{document}
